% ChargeShield-FL --- DSN 2027 --- 4. Experimental Methodology
% Scritta il 2026-09-23. Fonti: ESPERIMENTI.md (sez. 0 congelata il
% 2026-09-22 nella scheda risultati/decision_matrix_ACN_membership_DP.xlsx,
% regole, E-A...E-E), SISTEMA.md sez. 7 e 9, CanaryPositiveControl.md.
% Lo stato delle run nella Tabella tab:matrix e' quello del 2026-09-23 e va
% aggiornato dalle matrici, non dai log.

\section{Experimental Methodology}
\label{sec:method}

\subsection{Specification Fixed Before Reading the Results}
\label{sec:method-spec}

The guide of the project requires metrics, surfaces and thresholds to be
fixed before the confirmatory runs are read, not after. Table~\ref{tab:spec}
lists the decisions recorded on 2026-09-22, before any result of the
$\varepsilon$ sweep was inspected.

\begin{table}[t]
\centering
\caption{Evaluation specification fixed before the $\varepsilon$ sweep.}
\label{tab:spec}
\footnotesize
\begin{tabular}{@{}p{2.0cm}p{5.9cm}@{}}
\toprule
Field & Decision \\
\midrule
Primary attack metric & raw loss (loss threshold) on the global model; per-record count with permutation $z$ on the composed LiRA score; calibrated LiRA and shadow attack as secondary \\
Primary surface & per-client update (A1--A3) for the per-record analysis; global model for loss threshold and shadow attack \\
Operating FPR & 1\%, chance level $\mathrm{TPR} = \mathrm{FPR}$; 0.1\% only where the number of records supports it \\
Acceptable cost & final reconstruction loss within 3 times the no-DP reference \\
Equivalence margin & AUC within 0.02 of 0.5, two one-sided tests (TOST) \\
Replication unit & the seed; 5 seeds per cell; copies of a canary are not replicates \\
\bottomrule
\end{tabular}
\end{table}

\Susy{Guida sez. 8: il margine AUC 0.02 va motivato e non vale automaticamente
per il TPR; manca ancora un margine per il TPR e per la ricostruzione.}

\subsection{Attacks and Scores}
\label{sec:method-attacks}

All attacks derive their judgement from the reconstruction error of the target
record. The \emph{loss-threshold} attack scores $-\ell(x, f)$ on the global
model and uses no calibration. The \emph{shadow} attack scores the difference
between the loss of a shadow model, trained on half of the members, and the
loss of the target. \emph{LiRA} works on the per-client update: for each site
it trains 16 shadow models, retrained at every round with a warm start and
privatised like the client, fits a Gaussian log-likelihood ratio per record,
and sums the per-round scores into a \emph{composed score}. We report the AUC
and the TPR at 1\% FPR for each attack, and the reconstruction loss of the
final global model on the fixed natural hold-out, both as an absolute value and
as a ratio to the no-DP reference; since the no-DP loss is small, the ratio
exaggerates changes near zero and is read together with the absolute value.

\subsection{Per-Record Analysis}
\label{sec:method-perrecord}

Aggregate metrics can hide records that are recognised consistently. Within
each seed we convert the composed score of every member session into a
percentile. A session is \emph{flagged} if it is a member in at least two seeds,
with mean percentile at least 90 and minimum at least 75. The chance level is
obtained by permuting the scores within each seed, 200 times, and $z$ is the
distance of the observed count from the permutation mean in units of its
standard deviation. $z$ measures how far the count is from chance, not how
severe the exposure of any record is. Because the calibration of LiRA
degenerates in our regime (Section~\ref{sec:results-canary}), the composed
score is in practice a monotone function of the loss, and the analysis is a
ranking by loss; we describe it as such.

\subsection{Controls}
\label{sec:method-controls}

\paragraph{Canary positive control} On one client, $k = 20$ real sessions are
chosen as member templates and duplicated 30 times in the training data, and
$k = 20$ further sessions, drawn from the same pool, are moved to the
hold-out as non-member templates. The canary regime of
Section~\ref{sec:system-model} is used. Because the copies of a template are
identical, the AUC rests on $20 \times 20$ distinct pairs, not on the number
of copies.

\paragraph{Role exchange} Each run is repeated with the roles of the two sets
exchanged. A signal due to intrinsic difficulty is antisymmetric (the two AUCs
sum to one); a membership signal is symmetric (both exceed $0.5$). The exchange
is exact only when the two sets have the same size: an earlier configuration
with 5 member and 20 non-member templates drew disjoint sets in the two arms
and is not used.

\paragraph{Untrained baseline} The same canaries are scored on 20 randomly
initialised networks per arm and seed; the effect of training is the increase
over that baseline. The two arms' baselines sum to one by construction, so
their sum checks the exchange, not the data.

\paragraph{Clipping without noise} The guide foresees a clip-only arm (B1)
that separates the effect of clipping from that of noise. The code has no
clip-only setting yet; until it exists, the effect of client-level DP is
attributed to clipping and noise together. \Susy{segnalazione 38.}

\subsection{Experimental Matrix}
\label{sec:method-matrix}

Table~\ref{tab:matrix} links every experiment to its research question and
shows its status. Seeds are 42, 123, 456, 789 and 1234 in every cell.

\begin{table*}[t]
\centering
\caption{Experiments, research questions and status on 2026-09-23.}
\label{tab:matrix}
\footnotesize
\begin{tabular}{@{}p{1.7cm}p{1.2cm}p{4.4cm}p{4.2cm}p{4.6cm}@{}}
\toprule
ID & RQ & Factor varied & Held fixed & Status \\
\midrule
Campaign & RQ1 & no DP; client-level $\varepsilon \in \{0.1, 0.5, 1\}$ at A1, A2, A3; pilot $\varepsilon \in \{8, 16\}$ at A1 & data, split procedure, model, FedProx $\mu=0.01$, 10 rounds & completed; pilot cells have one seed (Section~\ref{sec:results-nodp}) \\
E-A & RQ1 & client-level $\varepsilon \in \{2, 4, 8, 16\}$, \texttt{dp-fedavg} (A1) & as above & running \pending{$\varepsilon=2$, seed 789: attacks failed on an environment error, to be re-run} \\
B1 & RQ1 & clipping without noise & as above & not executable (no clip-only setting) \\
E-B & RQ1 & record-level DP-SGD, noise multiplier $\in \{0.5, 1, 2\}$ & as above, client-level DP off & pending accountant and grouping fixes \\
E-C & RQ1 & per-record analysis of all new cells & --- & pending unification of the percentile computation \\
E-E & RQ3 & $\mu = 0$ against $\mu = 0.01$, without DP, then at the E-A operating point & data, split, seeds, rounds & pending configuration \\
E-D & RQ2 & IID against per-site partition, without DP, then at the E-A operating point & population, model, seeds & no-DP arm running \\
Canary, 2nd site & validity & balanced canary control on Caltech & canary protocol & pending \\
NVFlare & validity & deployment at the E-A operating point & analysis code & after E-A \\
\bottomrule
\end{tabular}
\end{table*}

\subsection{Statistics and Reading Rules}
\label{sec:method-stats}

The seed is the unit of replication and runs are paired by seed across cells.
Near-chance results are assessed with two one-sided tests against the margin of
Table~\ref{tab:spec}~\cite{schuirmann1987tost}; an interval that includes the
chance level is not by itself evidence of equivalence. With five paired seeds
an exact Wilcoxon signed-rank test cannot reach two-sided significance at
$0.05$ (its minimum $p$ is $0.0625$), so we report per-seed values rather than
significance claims for cell differences. For the canary we average the two
arms within each seed, which cancels the ceiling effect of low baselines, and
test the five per-seed means with a paired $t$ test and a one-sided sign test.

The conclusions allowed by each outcome of the $\varepsilon$ sweep were fixed
in advance (Table~\ref{tab:reading}).

\begin{table}[t]
\centering
\caption{Reading rules fixed before the $\varepsilon$ sweep (E-A).}
\label{tab:reading}
\footnotesize
\begin{tabular}{@{}p{2.7cm}p{2.7cm}p{2.3cm}@{}}
\toprule
Outcome & Allowed conclusion & Next step \\
\midrule
in a cell of acceptable cost the per-record excess falls below $z = 3$ & observed reduction of per-record exposure, for these attacks and this access & fix that $\varepsilon$ as operating point for E-E and E-D \\
excess remains at acceptable cost & client-level DP at that $\varepsilon$ does not reduce the measured exposure & report; E-B becomes decisive \\
cost above threshold in every cell & no useful operating point for client-level DP in this regime & answer to RQ1 for this unit; go to E-B \\
inconsistent across seeds & the experiment does not quantify the effect & targeted extra seeds on that cell only \\
\bottomrule
\end{tabular}
\end{table}

\subsection{Realisation}
\label{sec:method-realisation}

Experiments run in a single-process simulation. The same federation also runs
as a five-container NVIDIA FLARE deployment orchestrated with Containerlab
(server, three site clients and an administration console, with mutual TLS);
each round exports the updates and global weights, and the membership analysis
is run offline with the same functions as in simulation. Known differences are
that the deployment normalises features per site on the client, clips in
absolute terms in the first round under \texttt{dp-fedavg}, has no no-DP
switch, and requires the split seed to be read from the job snapshot. Main
results come from the simulation; the deployment is used to check that the
same conclusions hold.
